\documentclass[11pt,letterpaper]{article}
\usepackage[T1]{fontenc}
\usepackage{lmodern}
\usepackage{microtype}
\usepackage[margin=1in]{geometry}
\linespread{1.0}
\setlength{\parskip}{0pt}
\setlength{\parindent}{1.5em}
\usepackage{amsmath,amssymb,amsthm}
\usepackage{mathrsfs}
\usepackage{graphicx}
\usepackage{xcolor}
\usepackage{tikz}
\usepackage{booktabs,tabularx,array,longtable}
\newcolumntype{L}[1]{>{\raggedright\arraybackslash}p{#1}}
\usepackage{listings}
\lstset{basicstyle=\ttfamily\footnotesize,breaklines=true,columns=fullflexible,frame=single,framesep=4pt,showstringspaces=false,captionpos=b}
\usepackage[hidelinks]{hyperref}
\usepackage{cleveref}
\hypersetup{
 pdftitle={Prime-Adjoining Correspondences},
 pdfauthor={Matthew Long; The YonedaAI Collaboration},
 pdfsubject={Part I of Adelic Polarization Conjecture},
 pdfkeywords={semilocal adeles, Sonin spaces, Hochschild homology, Weil explicit formula}
}
\newtheorem{theorem}{Theorem}[section]
\newtheorem{proposition}[theorem]{Proposition}
\newtheorem{lemma}[theorem]{Lemma}
\newtheorem{corollary}[theorem]{Corollary}
\theoremstyle{definition}
\newtheorem{definition}[theorem]{Definition}
\theoremstyle{remark}
\newtheorem{remark}[theorem]{Remark}
\newcommand{\R}{\mathbb R}
\newcommand{\C}{\mathbb C}
\newcommand{\Q}{\mathbb Q}
\newcommand{\Z}{\mathbb Z}
\newcommand{\T}{\mathcal T}
\newcommand{\supp}{\operatorname{supp}}
\newcommand{\Rea}{\operatorname{Re}}
\newcommand{\Ima}{\operatorname{Im}}
\newcommand{\ind}{\mathbf1}
\newcommand{\Sch}{\mathcal S}
\newcommand{\Son}{\mathbf S}
\newcommand{\Bcal}{\mathcal B}
\usepackage{everypage}
\definecolor{grokgray}{RGB}{110,110,110}
\AddEverypageHook{%
  \ifnum\value{page}=1
    \begin{tikzpicture}[remember picture,overlay]
      \node[rotate=90,anchor=south,font=\footnotesize\sffamily\color{grokgray},inner sep=0pt]
        at ([xshift=30pt,yshift=0.5\paperheight]current page.south west)
        {GrokRxiv:2026.09.local-correspondence\quad [\,math.NT\,]\quad 08 Sep 2026};
    \end{tikzpicture}
  \fi
}
\title{Prime-Adjoining Correspondences\\[4pt]
  {\normalsize Part I of \emph{Adelic Polarization Conjecture}}}
\author{Matthew Long\\
  \small The YonedaAI Collaboration, YonedaAI Research Collective\\
  \small Chicago, IL\\
  \small \texttt{matthew@yonedaai.com}}
\date{September 2026}
\begin{document}
\maketitle
\begin{abstract}
A prime-adjoining map between semilocal Sonin spaces preserves Fourier duality but changes the inherited norm. We examine whether its local construction supplies a correspondence compatible with the arithmetic terms in Weil's explicit formula. Starting from the unit-invariant local Sonin vector, we derive the transition on logarithmic coordinates, its bounded inverse, its mixed pairing with the unit-ball map, and the two-prime compatibility square. The derivative of the logarithm of the positive metric operator has an absolutely convergent translation expansion whose distribution kernel contains every prime power with the required coefficient. This operator gives the local Weil term, whereas the raw norm defect has only three translation components. A separate algebraic obstruction concerns the direct tensor construction: no nonzero unit-invariant local Sonin vector defines a multiplicative coefficient map into the larger crossed-product algebra. The proof uses the incompatibility between the two shell values and idempotence. A weighted sum of shell embeddings instead gives Hochschild chain maps with commuting prime transitions, but its higher local tensor factors fail componentwise Fourier symmetry. Exact support indexing identifies when a finite-place formula equals the full Weil distribution. These results distinguish the analytic transition and corrected algebraic chain map from a still unconstructed primitive geometric pairing. The obstruction excludes a specified algebra-homomorphism ansatz; it does not exclude correspondences implemented by bimodules or corrected complexes, and it establishes no global positivity statement.
\end{abstract}

\section{Introduction}

Adjoining a finite prime to a semilocal adelic space changes a local Euler
factor and a Hilbert norm at the same time. The relation is explicit in the
Sonin-space construction of Connes, Consani and Moscovici
\cite{CCM}. For a prime $p$, the transition becomes multiplication by
$1-p^{-1/2-it}$ after passage to logarithmic coordinates and Fourier
transformation. Its modulus is bounded above and below by positive
constants. The resulting isomorphism therefore preserves the topology,
but it does not preserve the inherited inner product.

The arithmetic contribution of $p$ has a different form. Weil's explicit
formula involves every positive power of $p$ meeting the test support,
with coefficient $(\log p)p^{-k/2}$. A norm defect contains only the
constant coefficient and the first positive and negative translations.
The distinction determines which analytic construction could enter an
arithmetic comparison. The logarithm of the metric operator, followed by
differentiation with respect to the real parameter $\sigma$, produces the required infinite
translation series. The raw metric operator does not.

A further distinction appears before any question about positivity.
The local Sonin vector is a function supported on two $p$-adic shells.
Tensoring with that vector is a useful linear map, but it is not
multiplicative for pointwise multiplication. This obstruction survives
every nonzero scalar normalization of the unit-invariant local Sonin
vector. It prevents a direct interpretation of this tensor operation
as an algebra homomorphism between the natural semilocal crossed products.
It leaves other notions of correspondence available.

The geometric setting is substantial. Connes describes the semilocal
crossed products as a sheaf over $\operatorname{Spec}\Z$ and relates the
function spaces in the trace formula to Hochschild homology
\cite[Section 7.3, Theorem 7.2]{Connes2026}. Earlier work constructs
restriction-map cokernels and a descended trace pairing on the adele
class space \cite[Sections 4 and 5]{CCMar}. These results motivate a
primitive pairing, but they do not establish the required nonpositivity
of that geometric pairing $I$. The desired Weil form $B_W=-I$ would
be nonnegative. A passage from an analytic
map to a geometric correspondence must specify its action on the
relevant complex and its interaction with duality.

Our central algebraic result is
\cref{thm:local-correspondence:tensor-obstruction}. It rejects the direct
tensor algebra-homomorphism ansatz using only local fields and the actual
Sonin generator. A weighted sum of the two shell embeddings gives
an algebraic Hochschild chain map extending the degree-zero tensor
map, with commuting prime transitions. Its local factors fail
componentwise Fourier symmetry in every positive degree. The analytic
results surrounding these constructions give the exact maps
and distributions that a replacement would have to recover. Their proofs
also establish the signs and support conventions needed for comparisons
with a proposed intersection pairing. The standard local Fourier
calculation and the established Sonin stability theorem are attributed
to their sources; no novelty claim is made for them or for the elementary
Euler-factor expansion.

Part II develops different obstructions, concerning normalization on an
entire carrier and the sign of orthogonal primitive extensions
\cite{PartII}. The present argument takes place before those choices:
it tests multiplication on the proposed coefficient map itself. A
primitive polarization, a complete arithmetic comparison, and a global
spectral realization remain open construction problems.

\section{The complex test algebra}

\begin{definition}\label{def:local-correspondence:test-algebra}
Let $\T=C_c^\infty(\R,\C)$. For $f,g\in\T$, define
\begin{align}
 (f*g)(x)&=\int_\R f(u)g(x-u)\,du, &
 f^*(x)&=\overline{f(-x)},\label{eq:test-operations}\\
 F_f(z)&=\int_\R f(x)e^{zx}\,dx, &
 \widehat f(t)&=F_f(-it).\label{eq:test-transforms}
\end{align}
The Fourier inverse has the normalization
\begin{equation}\label{eq:fourier-inverse}
 f(x)=\frac1{2\pi}\int_\R\widehat f(t)e^{itx}\,dt.
\end{equation}
Operator-norm calculations use the unitary normalization
$(2\pi)^{-1/2}\widehat f$ on $L^2(\R)$; the transform
$\widehat f$ in the explicit formula retains
\eqref{eq:fourier-inverse}. All Hermitian forms in this paper
are linear in the first variable.
\end{definition}

The use of complex tests is necessary for the sesquilinear comparisons.
No evenness or reality condition is imposed on $f$. Even functions on
the real local field will occur in the Sonin realization, but their
logarithmic-coordinate functions need not be even. These are different
uses of the real variable, and identifying them would unnecessarily
restrict the test algebra.

\begin{lemma}\label{lem:test-identities}
For $f,g\in\T$, the function $F_f$ is entire and
\begin{align}
 F_{f*g}(z)&=F_f(z)F_g(z),\label{eq:laplace-convolution}\\
 F_{f^*}(z)&=\overline{F_f(-\overline z)},\label{eq:laplace-star}\\
 \widehat{f*g^*}(t)&=\widehat f(t)\overline{\widehat g(t)}.
 \label{eq:fourier-correlation}
\end{align}
If $\supp f\subset[-a,a]$, then for every $N\ge0$ and every $b>0$
there is a constant $C_{f,N,b}$ such that
\begin{equation}\label{eq:strip-decay}
 |F_f(\sigma+it)|\le C_{f,N,b}(1+|t|)^{-N}
 \qquad (|\sigma|\le b).
\end{equation}
\end{lemma}

\begin{proof}
Differentiation under the integral is valid on every compact subset of
$\C$, because $x$ lies in a fixed compact interval. The change of
variables $x=u+v$ proves \eqref{eq:laplace-convolution} by Fubini's
theorem. Replacing $x$ by $-x$ in the involution integral proves
\eqref{eq:laplace-star}. On the imaginary axis the latter identity is
complex conjugation, which gives \eqref{eq:fourier-correlation}.

For the decay estimate, integrate the Fourier transform of
$e^{\sigma x}f(x)$ by parts $N$ times. Its $N$th derivative has a
uniform $L^1$ bound for $|\sigma|\le b$, since the derivative is a finite
sum of terms $\sigma^j e^{\sigma x}f^{(N-j)}(x)$. The boundary terms
vanish. This proves the stated bound for $|t|\ge1$; a direct integral
bound handles the remaining compact interval.
\end{proof}

For $h=f*g^*$ the real-space formula is
\begin{equation}\label{eq:correlation}
 h(x)=\int_\R f(u)\overline{g(u-x)}\,du.
\end{equation}
Consequently $h(0)=\langle f,g\rangle_{L^2}$, and $h=h^*$ when $g=f$.
The Fourier transform of an autocorrelation is nonnegative. Its values
at nonzero real translations can nevertheless have either sign, even
for real tests. This is the elementary source of sign changes in local
prime contributions.

The multiplicative test corresponding to $h$ is
\begin{equation}\label{eq:multiplicative-test}
 \varphi_h(y)=y^{-1/2}h(\log y),\qquad y>0.
\end{equation}
With Mellin transform $\widetilde\varphi(s)=\int_0^\infty
\varphi(y)y^{s-1}\,dy$, substitution $y=e^x$ gives
\begin{equation}\label{eq:mellin-shift}
 \widetilde\varphi_h(s)=F_h(s-1/2).
\end{equation}
The shift is responsible for both $p^{-k/2}$ in the local terms and the
reflection $z\mapsto-\overline z$ in \eqref{eq:laplace-star}.
The conventions agree with Appendix B of Connes and Consani
\cite{CC2020} after this change of variables.

\section{The complete arithmetic distribution}

\begin{definition}\label{def:local-correspondence:weil}
For $h\in\T$, put
\begin{align}
 P(h)&=F_h(1/2)+F_h(-1/2),\label{eq:poles}\\
 W_p(h)&=(\log p)\sum_{k\ge1}p^{-k/2}
          \bigl(h(k\log p)+h(-k\log p)\bigr),\label{eq:local-weil}\\
 A(h)&=(\log(4\pi)+\gamma)h(0)\nonumber\\
 &\quad+\int_0^\infty
 \frac{e^{-u/2}(h(u)+h(-u))-2e^{-u}h(0)}{1-e^{-2u}}\,du.
 \label{eq:archimedean}
\end{align}
Here $\gamma$ is Euler's constant. The completed Weil distribution is
\begin{equation}\label{eq:weil-complete}
 W(h)=P(h)-A(h)-\sum_pW_p(h),\qquad
 B_W(f,g)=W(f*g^*).
\end{equation}
\end{definition}

Every finite-place sum in \eqref{eq:weil-complete} is finite on a fixed
compactly supported test. The apparently singular archimedean integral
also converges. Near zero,
$h(u)+h(-u)=2h(0)+O(u^2)$, so its numerator is $u h(0)+O(u^2)$.
The denominator is $2u+O(u^2)$. Beyond the support of $h$, only the
exponentially decaying subtraction term remains. Removing that term
before integration would destroy the cancellation at zero.

\begin{proposition}\label{prop:explicit-formula}
Let $\rho$ run over the nontrivial zeros of $\zeta$, with multiplicity.
Then
\begin{equation}\label{eq:zero-sum}
 W(h)=\sum_\rho F_h(\rho-1/2),\qquad h\in\T.
\end{equation}
The series converges absolutely. In particular,
\begin{equation}\label{eq:mixed-zero-sum}
 B_W(f,g)=\sum_\rho F_f(\rho-1/2)
      \overline{F_g(1/2-\overline\rho)}.
\end{equation}
\end{proposition}

\begin{proof}
Apply the explicit formula in \cite[Appendix B, equations (147) to
(150)]{CC2020} to \eqref{eq:multiplicative-test}. The two pole integrals
are $F_h(1/2)$ and $F_h(-1/2)$. The two values at $p^k$ become
$p^{-k/2}h(k\log p)$ and $p^{-k/2}h(-k\log p)$.
In the archimedean integral, $y=e^u$ changes
$dy/(y-y^{-1})$ into $du/(1-e^{-2u})$. This gives
\eqref{eq:archimedean} with the displayed signs.

The nontrivial zeros lie in $0<\Rea\rho<1$. The classical zero-count
bound is $O(T\log(T+2))$ up to height $T$; the standard counting formula
is recalled in \cite[Section 2]{Connes2026}. Taking $N>2$ in
\eqref{eq:strip-decay} and grouping zeros by dyadic height proves
absolute convergence. Finally, apply
\eqref{eq:laplace-convolution} and \eqref{eq:laplace-star} to $f*g^*$.
\end{proof}

The quadratic summand in \eqref{eq:mixed-zero-sum} is an absolute square
when $\Rea\rho=1/2$. Without that condition, its two factors are values
at different points. Absolute convergence therefore supplies no sign.
The reflection symmetry of the zero set does give a Hermitian form;
it does not make every summand nonnegative.

The pole form also consists of mixed products:
\begin{equation}\label{eq:pole-mixed}
 P(f*g^*)=F_f(1/2)\overline{F_g(-1/2)}
          +F_f(-1/2)\overline{F_g(1/2)}.
\end{equation}
The two products vanish if the first test has both pole moments zero.
This motivates the moment-killing subspace
\begin{equation}\label{eq:primitive-tests}
 \T_0=\{f\in\T:F_f(1/2)=F_f(-1/2)=0\}.
\end{equation}
On $\T_0$, the pole contribution disappears from the quadratic
explicit formula. Appendix C, Proposition 1 of \cite{CC2020},
with its finite exceptional set equal to $\{0,1\}$ in Mellin
coordinates, states that $B_W(f,f)\ge0$ for all $f\in\T_0$
is equivalent to RH. This precise version of the Weil criterion
justifies the two moment conditions. A further restriction to
even tests or an arbitrary smaller subspace requires its own
argument.

\section{The local shell vector}

Normalize additive Haar measure on $\Q_p$ by
$\operatorname{vol}(\Z_p)=1$, and take
$e_p(x)=\exp(2\pi i\{x\}_p)$, where $\{x\}_p$ is the standard
rational $p$-adic fractional part in $[0,1)$. This character has
conductor $\Z_p$. Let $\mathscr F_pu(y)=\int_{\Q_p}u(x)e_p(xy)\,dx$
on $L^1(\Q_p)\cap L^2(\Q_p)$, extended unitarily to
$L^2(\Q_p)$. For $n\in\Z$ write
\begin{equation}\label{eq:balls-shells}
 B_n=\{x:|x|_p\le p^n\},\qquad
 \epsilon_n=\ind_{B_n}-\ind_{B_{n-1}}.
\end{equation}
Thus $\epsilon_n$ is a shell indicator. Its volume is
$p^n-p^{n-1}$.

\begin{definition}\label{def:local-correspondence:shell}
The local Sonin vector is
\begin{equation}\label{eq:sigma}
 \sigma_p=\epsilon_0-p^{-1}\epsilon_1
 =(1+p^{-1})\ind_{B_0}-\ind_{B_{-1}}-p^{-1}\ind_{B_1}.
\end{equation}
The radial local Sonin space consists of $\Z_p^\times$-invariant
$u\in L^2(\Q_p)$ for which both $u$ and $\mathscr F_pu$ vanish almost
everywhere on $B_{-1}$.
\end{definition}

The following result is the local statement proved in
\cite[Section 4.5]{CCM}. A direct proof fixes the shell convention and
will be used in the algebraic obstruction.

\begin{proposition}\label{prop:local-sonin}
The radial local Sonin space is $\C\sigma_p$. Its generator satisfies
\begin{equation}\label{eq:sigma-properties}
 \mathscr F_p\sigma_p=\sigma_p,\qquad
 \int_{\Q_p}\sigma_p(x)\,dx=0,\qquad
 \|\sigma_p\|_2^2=1-p^{-2}.
\end{equation}
\end{proposition}

\begin{proof}
The Fourier transform of a ball is
\begin{equation}\label{eq:ball-fourier}
 \mathscr F_p\ind_{B_n}=p^n\ind_{B_{-n}}.
\end{equation}
Indeed, the character is trivial on the ball precisely on its
annihilator, and otherwise translation by an element of the ball forces
the integral to vanish. Apply \eqref{eq:ball-fourier} to the last
expression in \eqref{eq:sigma}. The coefficients of the three balls
remain unchanged. The shell description gives both the integral and
the squared norm:
\begin{align}
 \int\sigma_p&=(1-p^{-1})-p^{-1}(p-1)=0,\nonumber\\
 \|\sigma_p\|_2^2&=(1-p^{-1})+p^{-2}(p-1)=1-p^{-2}.
 \label{eq:sigma-norm-derivation}
\end{align}

For uniqueness, write a radial $u$ as a constant $a_n$ on each shell.
The spatial gap makes $a_n=0$ for $n<0$. For every integer $k\ge1$,
the Fourier gap and Plancherel give
\begin{equation}\label{eq:radial-moment}
 0=\langle\mathscr F_pu,\ind_{B_{-k}}\rangle
   =p^{-k}\int_{B_k}u(x)\,dx.
\end{equation}
These finite-ball integrals exist by the Cauchy-Schwarz inequality. Subtract the
identities for $k$ and $k-1$, when $k\ge2$, to obtain $a_k=0$.
The remaining identity for $k=1$ reads
$a_0(1-p^{-1})+a_1(p-1)=0$. It gives $a_1=-a_0/p$, as required.
\end{proof}

Replacing shells by balls would give a different function.
For example, $\ind_{B_0}-p^{-1}\ind_{B_1}$ equals $1-p^{-1}$ on
$B_{-1}$. It does not have the spatial Sonin gap. The subtraction
$\ind_{B_{-1}}$ in \eqref{eq:sigma} is essential, even though the
orbit-sum formula later has only two nonzero terms.

For $p=2$, the values of $\sigma_p$ on the shells $|x|_2=1$ and
$|x|_2=2$ are $1$ and $-1/2$, and its squared norm is $3/4$.
The negative shell value is compatible with Fourier self-duality.
It will be incompatible with idempotence for pointwise multiplication.
The two requirements concern different operations.

\section{Semilocal transitions}

Let $S$ be a finite set of places containing infinity, and put
$S_f=S\setminus\{\infty\}$. The semilocal objects are
\begin{equation}\label{eq:semilocal-objects}
 \mathbb A_S=\R\times\prod_{p\in S_f}\Q_p,
 \qquad \Gamma_S=\{\pm\prod_{p\in S_f}p^{n_p}:n_p\in\Z\},
 \qquad Y_S=\mathbb A_S/\Gamma_S.
\end{equation}
Let $K_S$ be the maximal compact subgroup of the semilocal idele class
group. We use the Hilbert realization
$H_S=L^2(Y_S)^{K_S}$ and its Fourier transform as constructed in
\cite[Section 4.1]{CCM}. This notation refers to that realization,
rather than to an unspecified quotient measure on a potentially
non-Hausdorff orbit space.

At the real place, $H_\infty=L^2(\R)^{\rm ev}$. Its multiplicative
coordinate map may be normalized as
\begin{equation}\label{eq:w-infinity}
 (w_\infty v)(y)=\sqrt2\,y^{1/2}v(y),\qquad y>0.
\end{equation}
It is unitary for measure $d^*y=dy/y$. Composing with $y=e^x$ gives
a unitary $J:H_\infty\to L^2(\R,dx)$,
$Jv(x)=\sqrt2e^{x/2}v(e^x)$. The factor $\sqrt2$ accounts for the
two equal halves of an even function. A common harmless Fourier
normalization factor is included whenever a unitary Fourier transform
is used.

\begin{definition}\label{def:local-correspondence:theta}
For $v\in H_\infty$, let
\begin{equation}\label{eq:theta-eta}
 \theta_Sv=[(\bigotimes_{p\in S_f}\sigma_p)\otimes v],\qquad
 \eta_Sv=[(\bigotimes_{p\in S_f}\ind_{\Z_p})\otimes v].
\end{equation}
Brackets denote the semilocal class realization of \cite{CCM}, initially
on its natural dense space and then by bounded extension. For
$S\subset T$, define the analytic transition
\begin{equation}\label{eq:relative-theta}
 \theta_{S,T}=\theta_T\theta_S^{-1}:H_S\longrightarrow H_T.
\end{equation}
\end{definition}

The one-prime orbit sum explains the operator without invoking a
geometric pairing. Evaluate at a representative whose finite coordinate
is $1$. Only the exponents $0$ and $-1$ of the $p$-orbit meet the two
shells. Up to the common normalization in \eqref{eq:w-infinity}, the
result is
\begin{equation}\label{eq:orbit-sum}
 y^{1/2}\bigl(v(y)-p^{-1}v(y/p)\bigr)
 =g(y)-p^{-1/2}g(y/p),\qquad g(y)=y^{1/2}v(y).
\end{equation}
This is the calculation in \cite[Section 4.6, equation (57)]{CCM}.
After $y=e^x$, it is a difference of two translations.

The corresponding unit-ball orbit sum contains all nonnegative
exponents. It gives
\begin{equation}\label{eq:eta-orbits}
 g(y)\longmapsto\sum_{k\ge0}p^{-k/2}g(p^ky).
\end{equation}
The two maps therefore translate in opposite directions before
Fourier transformation. This accounts for the complex conjugate in
their mixed pairing.

For $\lambda>0$, the real Sonin space is the subspace of even square
integrable functions whose spatial and real Fourier transforms vanish
on $(-\lambda,\lambda)$. The semilocal definition uses the corresponding
norm gap on $Y_S$. Theorem 4.6 of \cite{CCM} states that $\theta_S$
is a bounded topological isomorphism between these Sonin spaces.
The theorem includes surjectivity onto the gap-defined space; it is
stronger than the elementary multiplier bound alone. The Fourier
intertwining follows from $\mathscr F_p\sigma_p=\sigma_p$ and the
real factor in the tensor construction.

\section{The bounded translation operator}

Fix $p$, and abbreviate $L=\log p$ and $r=p^{-1/2}$. For real $b$,
let $(T_bu)(x)=u(x-b)$ on $L^2(\R,dx)$. Then
$T_b^*=T_{-b}$ and $T_bT_c=T_{b+c}$. Define
\begin{equation}\label{eq:Ap}
 A_p=I-rT_L,\qquad m_p(t)=1-re^{-itL}.
\end{equation}
Our Fourier transform sends $A_p$ to multiplication by $m_p$.

\begin{theorem}\label{thm:local-correspondence:inverse}
The operator $A_p$ is bounded and invertible, with
\begin{align}
 A_p^{-1}&=\sum_{k\ge0}r^kT_{kL},\label{eq:Ap-inverse}\\
 (1-r)\|u\|_2&\le\|A_pu\|_2\le(1+r)\|u\|_2,\label{eq:Ap-bounds}\\
 \|A_p\|&=1+r,\qquad \|A_p^{-1}\|=(1-r)^{-1}.
 \label{eq:Ap-exact-norms}
\end{align}
The series in \eqref{eq:Ap-inverse} converges in operator norm.
Its remainder after $k=N$ has norm at most $r^{N+1}/(1-r)$.
\end{theorem}

\begin{proof}
Since $\|rT_L\|=r<1$, the geometric series is absolutely convergent
in the Banach algebra of bounded operators. Multiplication of its
partial sum by $I-rT_L$ leaves $I-r^{N+1}T_{(N+1)L}$. The remainder
bound proves both inverse identities. The triangle and reverse
triangle inequalities give \eqref{eq:Ap-bounds}.

Under the unitary transform $(2\pi)^{-1/2}\widehat{\phantom f}$,
the operator norms are the essential suprema of $|m_p|$ and
$|m_p|^{-1}$. The continuous periodic function $|m_p|$ takes its
minimum $1-r$ at $t=0$ and its maximum $1+r$ at $t=\pi/L$.
Every neighborhood of either point has positive measure. The essential
extrema therefore equal the stated extrema.
\end{proof}

The inverse is an ambient Hilbert-space inverse. It usually fails to
preserve compact support. To see the failure exactly, choose a nonzero
$u\in\T$ supported in an interval of length less than $L$. The
translates $T_{kL}u$ have disjoint supports, and on the $k$th translated
interval the inverse equals $r^kT_{kL}u$. Infinitely many of these
pieces are nonzero. The inverse remains square integrable because the
sum of their squared norms is a geometric series.

By contrast, $A_p$ preserves smooth compact support and obeys
\begin{equation}\label{eq:support-Ap}
 \supp(A_pu)\subset\supp u\cup(\supp u+L).
\end{equation}
It acts continuously on every Sobolev space, with the same bounds,
because its translation terms commute with weak derivatives.
These properties are useful for local analytic transitions but do
not make $A_p$ an automorphism of the fixed compact test algebra.

For a nonzero compactly supported smooth function $u$ and real $b$,
the commutator with multiplication by $x$ is
\begin{equation}\label{eq:position-commutator}
 [x,T_b]u=bT_bu,\qquad [x,A_p]u=-rLT_Lu.
\end{equation}
Thus the coefficient $\log p$ can already be read from the displacement
in the local operator. The logarithmic derivative below recovers it
with all iterates of that displacement.

\section{Mixed duality and metric variation}

The unit-ball operator in logarithmic coordinates is
\begin{equation}\label{eq:Bp}
 B_p=(A_p^*)^{-1}=\sum_{k\ge0}r^kT_{-kL}.
\end{equation}
It has Fourier multiplier $1/\overline{m_p(t)}$. The orientation of
the translation in \eqref{eq:Bp} is fixed by the orbit sum
\eqref{eq:eta-orbits}; replacing it by $A_p^{-1}$ would change the
mixed pairing.

\begin{proposition}\label{prop:local-correspondence:mixed}
For all $u,v\in L^2(\R)$,
\begin{equation}\label{eq:mixed-pairing}
 \langle A_pu,B_pv\rangle=\langle u,v\rangle.
\end{equation}
For finite $S$, the semilocal maps satisfy
\begin{equation}\label{eq:semilocal-mixed}
 \langle\theta_Su,\eta_Sv\rangle_{H_S}
   =\langle u,v\rangle_{H_\infty}.
\end{equation}
\end{proposition}

\begin{proof}
Since $A_p^*B_p=I$, the adjoint identity proves
\eqref{eq:mixed-pairing}. Equivalently, its Fourier integrand contains
$m_p\overline{1/\overline{m_p}}=1$. The finite products over $S_f$
commute, and their factors cancel in the same way. The unitary
semilocal realization gives \eqref{eq:semilocal-mixed}, which is the
mixed duality identity of \cite[Section 4.7]{CCM}.
\end{proof}

The mixed identity pairs two different embeddings. Its specialization
to $u=v$ is positive, but it is not a self-pairing of $\theta_Su$.
The latter is controlled by the positive operator
\begin{equation}\label{eq:Gp}
 G_p=A_p^*A_p=(1+r^2)I-r(T_L+T_{-L}).
\end{equation}

\begin{proposition}\label{prop:local-correspondence:defect}
The defect $D_p=G_p-I$ satisfies
\begin{align}
 D_p&=r^2I-r(T_L+T_{-L}),\label{eq:Dp}\\
 \langle A_pu,A_pv\rangle-\langle u,v\rangle
   &=\langle D_pu,v\rangle.\label{eq:metric-defect}
\end{align}
Its Fourier multiplier is
\begin{equation}\label{eq:dp-symbol}
 d_p(t)=r^2-2r\cos(tL).
\end{equation}
The defect has both positive and negative quadratic values on the
ambient Hilbert space.
\end{proposition}

\begin{proof}
Expand $A_p^*A_p$ using $T_{-L}T_L=I$. The adjoint convention gives
\eqref{eq:metric-defect}. At $t=0$, $d_p=r^2-2r<0$; at
$t=\pi/L$, $d_p=r^2+2r>0$. Take nonzero Fourier transforms supported
in sufficiently small neighborhoods of the two points. Plancherel
then makes the corresponding quadratic integrals strictly negative
and strictly positive.
\end{proof}

The localized witnesses establish the ambient signs. Membership in a
fixed Sonin space requires simultaneous spatial and real Fourier
gaps, which the logarithmic-frequency localization used in the
proof does not preserve.

The defect has a simple real-space test. If $u$ is supported in an
interval of length less than $L$, then
$\langle T_Lu,u\rangle=\langle T_{-L}u,u\rangle=0$, and
\begin{equation}\label{eq:short-bump-defect}
 \langle D_pu,u\rangle=r^2\|u\|_2^2>0.
\end{equation}
This is already incompatible with identifying the defect with the
local Weil autocorrelation term: that arithmetic term is zero for
the same short bump. A larger support gives a stronger distinction
by detecting the missing second prime power.

\section{The common entire carrier}

Write the standard local factors as
\begin{equation}\label{eq:Euler-factors}
 L_\infty(s)=\pi^{-s/2}\Gamma(s/2),\qquad
 L_p(s)=(1-p^{-s})^{-1},\qquad
 E_S(t)=\prod_{v\in S}L_v(1/2+it).
\end{equation}
The dual Hardy-Titchmarsh realization of \cite[Section 4.4]{CCM}
uses the weighted space with density $|E_S(t)|^{-2}$. The Sonin
spaces map to a common space $\Bcal_\lambda$ of entire functions;
the boundary value $F(1/2+it)$ carries the norm
\begin{equation}\label{eq:carrier-norm}
 q_S(F,G)=\int_\R F(1/2+it)\overline{G(1/2+it)}
               |E_S(t)|^{-2}\,dt.
\end{equation}
The normalization is fixed consistently across $S$. Equations (59)
to (62) of \cite{CCM} identify $\theta_{S,T}$ with the identity on
this common entire carrier.

For $p\notin S$ the weight relation is
\begin{equation}\label{eq:weight-relation}
 |E_{S\cup\{p\}}(t)|^{-2}=|m_p(t)|^2|E_S(t)|^{-2}.
\end{equation}
Consequently, with $w_S(t)=|E_S(t)|^{-2}$,
\begin{align}
 q_{S\cup\{p\}}(F,G)-q_S(F,G)
 &=\int_\R F(1/2+it)\overline{G(1/2+it)}d_p(t)w_S(t)\,dt,
 \label{eq:carrier-defect}\\
 (1-r)^2q_S(F,F)&\le q_{S\cup\{p\}}(F,F)
                  \le(1+r)^2q_S(F,F).
 \label{eq:carrier-comparison}
\end{align}
The identity on the common entire carrier is a topological
isomorphism. The changing inner products in
\eqref{eq:carrier-defect} do not make it an isometry.

The weighted ambient spaces have positive, locally finite densities.
Compactly supported frequency functions therefore prove both signs
for \eqref{eq:carrier-defect} on those ambient spaces. They do not
prove both signs on $\Bcal_\lambda$. The possible restriction of a
multiplier to the entire carrier must be established separately.

There is also a distinction between two polar decompositions. The
identity between the weighted ambient spaces has a pointwise polar
normalization by $|m_p|^{-1}$. The identity between the restricted
Hilbert norms on $\Bcal_\lambda$ has an adjoint computed from the
restricted forms. That adjoint can involve a compression and need not
be multiplication by $|m_p|^2$. The entire-carrier obstruction in
Part II, Theorem 3.4 \cite{PartII}, concerns the first, explicitly pointwise
normalization. It does not identify the second polar decomposition
with that multiplier.

\section{The analytic two-prime square}

Let $p\ne q$ be primes. Use subscripts to distinguish
$L_p=\log p$, $r_p=p^{-1/2}$, and the analogous quantities for $q$.
Their translation operators commute on $L^2(\R)$.

\begin{proposition}\label{prop:local-correspondence:square}
The prime-adjoining operators satisfy
\begin{equation}\label{eq:two-prime-A}
 A_pA_q=A_qA_p
 =I-r_pT_{L_p}-r_qT_{L_q}+r_pr_qT_{L_p+L_q}.
\end{equation}
The relative semilocal maps obey
\begin{equation}\label{eq:two-prime-theta}
 \begin{aligned}
 &\theta_{S\cup\{p\},S\cup\{p,q\}}\theta_{S,S\cup\{p\}}\\
 &\qquad=\theta_{S\cup\{q\},S\cup\{p,q\}}\theta_{S,S\cup\{q\}}.
 \end{aligned}
\end{equation}
The metric operator is $G_{pq}=G_pG_q$, and its defect is
\begin{equation}\label{eq:two-prime-defect}
 G_{pq}-I=D_p+D_q+D_pD_q.
\end{equation}
\end{proposition}

\begin{proof}
Expand $(I-r_pT_{L_p})(I-r_qT_{L_q})$ and use the addition law
for translations. Formula \eqref{eq:two-prime-theta} also follows
directly from \eqref{eq:relative-theta}; in the local tensor
description it is the interchange of the $p$ and $q$ factors.
Every $A_p$ commutes with every $A_q$ and $A_q^*$, so
$(A_pA_q)^*(A_pA_q)=G_pG_q$. Substituting $G_p=I+D_p$ and
$G_q=I+D_q$ gives the final identity.
\end{proof}

The mixed defect in \eqref{eq:two-prime-defect} is explicit:
\begin{align}
 D_pD_q={}&r_p^2r_q^2I
 -r_p^2r_q(T_{L_q}+T_{-L_q})
 -r_q^2r_p(T_{L_p}+T_{-L_p})\nonumber\\
 &+r_pr_q\bigl(T_{L_p+L_q}+T_{L_p-L_q}
            +T_{-L_p+L_q}+T_{-L_p-L_q}\bigr).
 \label{eq:mixed-translations}
\end{align}
It includes translations by $\log(pq)$ and $\log(p/q)$.
The local prime sum in Weil's formula has no term indexed by $pq$
when $p\ne q$. Thus additive accumulation of raw metric defects
cannot reproduce the arithmetic prime decomposition.

The two orders of adjoining primes still agree. In the weighted
carrier, the defect from adjoining $q$ after $p$ has density
$d_q(1+d_p)w_S$; the other first step has density $d_pw_S$.
The sum is $(d_p+d_q+d_pd_q)w_S$. Reversing the order gives the
same expression. Commutativity controls the analytic square while
retaining a nonzero mixed correction.

The logarithm behaves differently. Since the positive operators
$G_p,G_q$ are commuting multiplication operators after Fourier
transformation,
\begin{equation}\label{eq:log-additivity}
 \log(G_pG_q)=\log G_p+\log G_q.
\end{equation}
No branch ambiguity occurs for the logarithms of positive invertible
operators. This exact additivity suggests where to seek an arithmetic
local term. It remains an identity of analytic operators and does
not construct a square of geometric primitive correspondences.

\section{Logarithmic variation}

For a real parameter $\sigma>0$, put $r_\sigma=p^{-\sigma}$ and
\begin{equation}\label{eq:sigma-family}
 A_{p,\sigma}=I-r_\sigma T_L,\qquad
 G_{p,\sigma}=A_{p,\sigma}^*A_{p,\sigma}.
\end{equation}
These are bounded invertible operators. The spectrum of
$G_{p,\sigma}$ lies in $[(1-r_\sigma)^2,(1+r_\sigma)^2]$,
so the self-adjoint logarithm is defined by continuous functional
calculus.

\begin{theorem}\label{thm:local-correspondence:logarithmic}
The logarithm and its derivative admit operator-norm convergent series
\begin{align}
 \log G_{p,\sigma}
 &=-\sum_{k\ge1}\frac{r_\sigma^k}{k}
                (T_{kL}+T_{-kL}),\label{eq:operator-log}\\
 K_{p,\sigma}:=\frac{d}{d\sigma}\log G_{p,\sigma}
 &=L\sum_{k\ge1}r_\sigma^k(T_{kL}+T_{-kL}).
 \label{eq:operator-log-derivative}
\end{align}
Differentiation is locally uniform for $\sigma>0$. For $f,g\in\T$,
\begin{equation}\label{eq:operator-local-Weil}
 \langle K_{p,1/2}f,g\rangle=W_p(f*g^*).
\end{equation}
Equivalently, for every $h\in\T$,
\begin{equation}\label{eq:scalar-local-Weil}
 W_p(h)=\frac1{2\pi}\int_\R\widehat h(t)
 \left.\frac{\partial}{\partial\sigma}
       \log|1-p^{-\sigma-it}|^2\right|_{\sigma=1/2}\,dt.
\end{equation}
\end{theorem}

\begin{proof}
For $|z|<1$, the analytic logarithm has power series
$\log(1-z)=-\sum_{k\ge1}z^k/k$. Substitute
$z=r_\sigma e^{-itL}$ and its conjugate, then add. The resulting
real function is the logarithm of
$|1-r_\sigma e^{-itL}|^2$. Its series is uniformly absolutely
convergent in $t$. Fourier multiplication gives
\eqref{eq:operator-log}, with absolute convergence bounded by
$2\sum_{k\ge1}r_\sigma^k/k$.

On $\sigma\ge\sigma_0>0$, the series of derivatives is bounded
in norm by $2L\sum_{k\ge1}p^{-k\sigma_0}$. The differentiated
series therefore converges uniformly on every compact positive
parameter interval and equals the derivative of the logarithm.
Its Fourier multiplier is
\begin{equation}\label{eq:log-derivative-symbol}
 k_{p,\sigma}(t)=2L\sum_{k\ge1}p^{-k\sigma}\cos(ktL).
\end{equation}

For $h=f*g^*$, \eqref{eq:correlation} gives
$\langle T_{kL}f,g\rangle=h(-kL)$ and
$\langle T_{-kL}f,g\rangle=h(kL)$. Substitution in
\eqref{eq:operator-log-derivative} proves
\eqref{eq:operator-local-Weil}. Alternatively,
$\widehat h\in L^1$ and the uniformly absolutely convergent
series \eqref{eq:log-derivative-symbol} permit integration term by
term. Fourier inversion gives
\begin{equation}\label{eq:cosine-inversion}
 \frac1{2\pi}\int_\R\widehat h(t)\cos(ktL)\,dt
     =\frac12\bigl(h(kL)+h(-kL)\bigr),
\end{equation}
which proves \eqref{eq:scalar-local-Weil} for arbitrary $h\in\T$.
\end{proof}

The finite measure
\begin{equation}\label{eq:kernel-measure}
 \mu_{p,\sigma}=L\sum_{k\ge1}p^{-k\sigma}
                     (\delta_{kL}+\delta_{-kL})
\end{equation}
is the convolution kernel of $K_{p,\sigma}$. Its total variation is
$2Lp^{-\sigma}/(1-p^{-\sigma})$. The measure is positive as a
measure, yet convolution by it is not a positive operator. Positivity
of a convolution operator is a condition on its Fourier multiplier;
positive coefficients in the real-space measure do not establish it.

Indeed, at $t=0$ the multiplier equals
$2Lr_\sigma/(1-r_\sigma)>0$. At $t=\pi/L$ it equals
$-2Lr_\sigma/(1+r_\sigma)<0$. Thus the operator derivative has both
signs on the ambient Hilbert space. Its exact norm is
\begin{equation}\label{eq:K-norm}
 \|K_{p,\sigma}\|=\frac{2Lp^{-\sigma}}{1-p^{-\sigma}}.
\end{equation}
The upper bound follows from the series, and concentration near
$t=0$ gives equality.

The resolvent form is also useful:
\begin{equation}\label{eq:K-resolvent}
 K_{p,\sigma}=L\left(r_\sigma T_L(I-r_\sigma T_L)^{-1}
             +r_\sigma T_{-L}(I-r_\sigma T_{-L})^{-1}\right).
\end{equation}
All factors commute, so this is also
$G_{p,\sigma}^{-1}(dG_{p,\sigma}/d\sigma)$. For a noncommuting
family, that last expression need not be the derivative of the
operator logarithm. The commuting translation structure is doing
specific work here.

\section{Prime powers and finite propagation}

The difference between $D_p$ and $K_{p,1/2}$ can be detected with
compact smooth functions. Choose a nonzero real $\psi$ supported
in $(-\varepsilon,\varepsilon)$, with $2\varepsilon<L$.
For $h=\psi*\psi^*$, all values $h(kL)$ with nonzero integer $k$
vanish. Hence
\begin{equation}\label{eq:short-bump-comparison}
 W_p(h)=0,\qquad \langle D_p\psi,\psi\rangle=r^2\|\psi\|_2^2.
\end{equation}
Even a scalar multiple of the raw defect cannot equal $W_p$ on all
autocorrelations, unless that scalar is zero; the latter possibility
is excluded by the next example.

Let $f=\psi+T_{2L}\psi$. The two pieces have disjoint support, and
the only nonzero correlations at nonzero multiples of $L$ occur
at $\pm2L$. Their values are $\|\psi\|_2^2$. Thus
\begin{equation}\label{eq:second-prime-power}
 W_p(f*f^*)=2Lr^2\|\psi\|_2^2,
 \qquad \langle D_pf,f\rangle=2r^2\|\psi\|_2^2.
\end{equation}
The arithmetic quantity records $p^2$. A putative formula retaining
only $h(\pm L)$ would return zero. Combining this example with
\eqref{eq:short-bump-comparison} rules out a universal proportionality
between the full local term and the raw norm defect.

More generally, for a finite coefficient sequence $c_0,\ldots,c_n$,
put $f=\sum_{j=0}^nc_jT_{jL}\psi$. The exact finite formula is
\begin{equation}\label{eq:finite-toeplitz}
 W_p(f*f^*)=L\|\psi\|_2^2
       \sum_{\substack{0\le i,j\le n\\i\ne j}}
       c_i\overline{c_j}r^{|i-j|}.
\end{equation}
For example, $f=\psi+T_L\psi$ gives a positive local value and
$f=\psi-T_L\psi$ gives a negative one. These two examples do not
kill the pole moments. The stronger moment-killed examples, including
the explicit prime $11$ witnesses, are established in Part II,
Theorem 6.3 \cite{PartII}.

The algebraic structure of \eqref{eq:finite-toeplitz} also explains
why prime powers cannot be discarded after a support enlargement.
Every pair of separated bumps at distance $kL$ contributes the
$k$th coefficient. A truncation retaining only $k=1$ loses all
interactions between nonadjacent packets. Smoothness does not remove
them, since their values are exact inner products of translated
copies of the same function.

At the distribution level, the first omitted coefficient is already
enough to distinguish a truncation. Choose $h\in\T$ supported near
$2L$, equal to $1$ at $2L$, with support excluding $0$, $\pm L$ and
all other integer multiples of $L$. Then $W_p(h)=Lr^2$, while every
distribution supported on $\{0,L,-L\}$ vanishes on $h$. This test
compares the distributions directly without factoring $h$ as an
autocorrelation.

\section{Support-indexed arithmetic}

\begin{definition}\label{def:local-correspondence:support}
For $a>0$, let $\T_a=\{f\in\T:\supp f\subset[-a,a]\}$.
A pair $(S,a)$ is support-admissible if $S$ is a finite place set
containing infinity and every prime $p\le e^{2a}$. Define
\begin{equation}\label{eq:finite-Weil}
 W_S=P-A-\sum_{p\in S_f}W_p.
\end{equation}
The pairs are ordered by $(S,a)\le(T,b)$ when $S\subset T$ and
$a\le b$.
\end{definition}

\begin{proposition}\label{prop:local-correspondence:support}
For $f,g\in\T_a$, the local contribution $W_p(f*g^*)$ can be
nonzero only if $p^k\le e^{2a}$ for some $k\ge1$.
At a support-admissible pair,
\begin{equation}\label{eq:support-equality}
 W_S(f*g^*)=W(f*g^*),\qquad f,g\in\T_a.
\end{equation}
If $p\notin S$ and $(S,a)$ is support-admissible, then
\begin{equation}\label{eq:new-prime-zero}
 W_p(f*g^*)=0\qquad(f,g\in\T_a).
\end{equation}
\end{proposition}

\begin{proof}
The support of $f*g^*$ is contained in
$\supp f-\supp g\subset[-2a,2a]$. Hence a value at
$\pm k\log p$ can be nonzero only if $k\log p\le2a$.
The stated exponential condition is equivalent. All potentially
contributing primes are in $S$, which proves
\eqref{eq:support-equality}. A newly adjoined prime has
$p>e^{2a}$, proving \eqref{eq:new-prime-zero}.
\end{proof}

The weak inequality in the support bound is a safe indexing
convention. Smooth functions supported in a closed interval vanish
at its endpoints, so some boundary contributions are automatically
zero. Including the boundary primes makes no error and avoids a
dependence on the particular support of each test.

For unequal supports $\supp f\subset[-a,a]$ and
$\supp g\subset[-b,b]$, the sharper sufficient condition is
$p^k\le e^{a+b}$. For translated intervals, the difference set
$\supp f-\supp g$ is the exact support constraint. The symmetric
family $\T_a$ is convenient because it is directed and covers all
compact smooth tests. It is not the smallest possible place set
for every individual pair.

As a concrete support window, take $a=\log2$. Then
$e^{2a}=4$, and the finite places $2$ and $3$ suffice. The indexed
prime powers are $2$, $3$, and $4$, though the value at the endpoint
$\log4=2a$ vanishes for these smooth autocorrelations. Taking a
slightly larger $a$ makes the $p^2=4$ coefficient available in the
interior. A prime $5$ adjoined before the support reaches
$\tfrac12\log5$ contributes nothing to the quadratic tests in that
window.

The support-admissible pairs form a directed set. Given
$(S,a)$ and $(T,b)$, choose $c=\max(a,b)$ and enlarge $S\cup T$
by the finitely many primes at most $e^{2c}$. The resulting pair
dominates both. Since each test pair appears in one such window,
the finite-place values in \eqref{eq:support-equality} stabilize
exactly. No infinite Euler product is needed for this distributional
stabilization.

For an intermediate set $S$ that is not support-admissible,
\begin{equation}\label{eq:finite-increment}
 W_{S\cup\{p\}}(f*g^*)-W_S(f*g^*)=-W_p(f*g^*).
\end{equation}
The formula is exact, but the intermediate distribution need not
have the positivity properties desired for a complete arithmetic
stage. In a proposed comparison $B_W=-I$, the corresponding
intersection increment would have sign $+W_p$. Boundary and degree
corrections must be included before asserting such an identity for
a geometric construction.

\section{The tensor algebra obstruction}

The semilocal crossed product contains a natural coefficient
algebra. Let $\Sch(\mathbb A_S)$ be the Bruhat-Schwartz functions,
with pointwise multiplication. For $\gamma\in\Gamma_S$, define
$\alpha_\gamma f(x)=f(\gamma^{-1}x)$. The algebraic crossed product
\begin{equation}\label{eq:crossed-product}
 \mathcal A_S=\Sch(\mathbb A_S)\rtimes_{\rm alg}\Gamma_S
\end{equation}
consists of finite sums $\sum_\gamma f_\gamma U_\gamma$, with
\begin{equation}\label{eq:crossed-product-multiplication}
 (fU_\gamma)(gU_\delta)
        =f\alpha_\gamma(g)U_{\gamma\delta}.
\end{equation}
This is the smooth groupoid algebra convention used in
\cite[Section 4, equations (4.3) to (4.6)]{CCMar}.

For $p\notin S$, the inclusion $\Gamma_S\subset\Gamma_{S\cup\{p\}}$
is fixed. Given a unit-invariant $v\in\Sch(\Q_p)$, define a linear
coefficient map
\begin{equation}\label{eq:tensor-map}
 \Phi_v:\mathcal A_S\longrightarrow\mathcal A_{S\cup\{p\}},
 \qquad \Phi_v(fU_\gamma)=(v\otimes f)U_\gamma.
\end{equation}
This map retains the old group element and tensors the coefficient
with the proposed local vector. It is the direct algebraic candidate
associated with the linear tensor operation in $\theta$.

\begin{lemma}\label{lem:tensor-multiplication}
For all $a,b\in\mathcal A_S$,
\begin{equation}\label{eq:tensor-defect}
 \Phi_v(a)\Phi_v(b)=\Phi_{v^2}(ab).
\end{equation}
If the coefficient algebra is nonzero, $\Phi_v$ is multiplicative
if and only if $v^2=v$.
\end{lemma}

\begin{proof}
Every $\gamma\in\Gamma_S$ is a $p$-adic unit, since $p\notin S$.
The function $v$ is unit-invariant, so
$\alpha_\gamma(v\otimes g)=v\otimes\alpha_\gamma g$.
Apply \eqref{eq:crossed-product-multiplication} to two elementary
terms; their local coefficient is $v^2$. Bilinearity gives
\eqref{eq:tensor-defect}. If $v^2=v$, this proves multiplicativity.

Conversely choose a nonzero real $f\in\Sch(\mathbb A_S)$.
Its square is nonzero. Multiplicativity on the two identity-group
elements $fU_1$ would imply
$(v^2-v)\otimes f^2=0$. Evaluate at a point where $f^2$ is nonzero
to obtain $v^2-v=0$.
\end{proof}

\begin{theorem}\label{thm:local-correspondence:tensor-obstruction}
No nonzero unit-invariant local Sonin vector $v$ makes the map
$\Phi_v$ in \eqref{eq:tensor-map} multiplicative. In particular,
the shell tensor construction with $v=\sigma_p$ does not define
an algebra homomorphism between these crossed products, and no
nonzero scalar normalization repairs it.
\end{theorem}

\begin{proof}
By Proposition~\ref{prop:local-sonin}, every such vector is $v=c\sigma_p$.
It is therefore a Bruhat-Schwartz function, so
\eqref{eq:tensor-map} is defined. If $\Phi_v$ were multiplicative,
Lemma~\ref{lem:tensor-multiplication} would imply $v^2=v$.
On the shell $|x|_p=1$, that identity is $c^2=c$.
Since $v\ne0$, it gives $c=1$. On the shell $|x|_p=p$,
the same identity would then require $p^{-2}=-p^{-1}$,
which is impossible for a positive prime. This contradiction proves
both assertions.
\end{proof}

The multiplicativity defect is supported on just one shell:
\begin{equation}\label{eq:sigma-square}
 \sigma_p^2-\sigma_p=(p^{-2}+p^{-1})\epsilon_1.
\end{equation}
It is nonzero before taking any completion, quotient, trace, or
primitive projection. At $p=2$, the tensor map sends a coefficient
to $-1/2$ times that coefficient on the outer shell; multiplying
two images gives $1/4$ times their product, while the image of
their product remains $-1/2$ times it.

The unit-ball vector offers a useful comparison. For
$v=\ind_{\Z_p}$, one has $v^2=v$, and
Lemma~\ref{lem:tensor-multiplication} gives a multiplicative map.
That vector is Fourier fixed, but it does not vanish on $B_{-1}$
and hence does not lie in the local Sonin space. The two local
vectors support different features: the ball vector permits the
coefficient algebra map, while the shell vector supplies the
Sonin gap and the analytic $\theta$ transition.

The theorem addresses an exact ansatz, not every arithmetic
correspondence. A bimodule need not arise from an algebra
homomorphism of the form \eqref{eq:tensor-map}; a chain-level
construction can contain corrections; a trace pairing can descend
through a quotient by a different map. The theorem also does not
assert that $\theta$ was intended in \cite{CCM} as an algebra
homomorphism. It tests that additional interpretation and rejects it.

The calculation is independent of the Riemann Hypothesis and of
the zero distribution. Its inputs are the two local gaps,
unit invariance, and pointwise multiplication in the coefficient
algebra. It gives a precise place where a proposed geometric
upgrade of the analytic transition must change its construction.

\section{A corrected Hochschild map}

The complex in this section uses uncompleted algebraic tensor
products over $\C$. Its boundary uses multiplication alone, so it
is defined for the nonunital algebra $\mathcal A_S$ without adding
a unit or using unit-dependent degeneracy maps. All homology
statements refer to the displayed algebraic complex.

The two shell indicators define algebra maps separately. Write
$e_0=\epsilon_0$ and $e_1=\epsilon_1$. They satisfy
\begin{equation}\label{eq:orthogonal-shells}
 e_0^2=e_0,\qquad e_1^2=e_1,\qquad e_0e_1=0.
\end{equation}
Each is invariant under the old group $\Gamma_S$, so
Lemma~\ref{lem:tensor-multiplication} gives algebra embeddings
$\Phi_{e_j}:\mathcal A_S\to\mathcal A_{S\cup\{p\}}$.
Injectivity follows by evaluating the local coefficient at any
point of the nonempty shell. Their images multiply to zero in
either order, because the old group elements preserve both
shells. They lie in the subalgebra with group support in
$\Gamma_S$.

The ambient target has additional group elements, and those
elements need not preserve the shells. For example, the new
generator $p$ changes the $p$-adic norm by $p^{-1}$.
Thus the shell indicators are not asserted to be central
in the full target crossed product. The Bruhat-Schwartz
coefficient algebra is also nonunital at the real place;
no global identity coefficient is inserted in the argument.

The degree-zero shell map is the weighted difference
\begin{equation}\label{eq:signed-embeddings}
 \Phi_{\sigma_p}=\Phi_{e_0}-p^{-1}\Phi_{e_1}.
\end{equation}
Although that combination is not multiplicative, it has a
natural extension as a map of algebraic Hochschild complexes.
The extension uses the same coefficient at every degree,
rather than the tensor powers of the combination in
\eqref{eq:signed-embeddings}.

For a complex algebra $\mathcal A$, set
$C_n(\mathcal A)=\mathcal A^{\otimes(n+1)}$ for $n\ge0$,
where tensors are algebraic over $\C$. The Hochschild boundary
on an elementary tensor is
\begin{align}
 b(a_0\otimes\cdots\otimes a_n)
 ={}&\sum_{j=0}^{n-1}(-1)^j
 a_0\otimes\cdots\otimes a_ja_{j+1}\otimes\cdots\otimes a_n
 \nonumber\\
 &+(-1)^n a_na_0\otimes a_1\otimes\cdots\otimes a_{n-1}.
 \label{eq:Hochschild-boundary}
\end{align}
The boundary from degree zero is zero. Associativity implies
$b^2=0$ by cancellation of adjacent multiplication terms.
These definitions do not require a unit.

\begin{proposition}\label{prop:corrected-chain-map}
For each $n\ge0$, define
\begin{equation}\label{eq:corrected-chain-map}
 \Psi_n=\Phi_{e_0}^{\otimes(n+1)}
                -p^{-1}\Phi_{e_1}^{\otimes(n+1)}.
\end{equation}
Then $b\Psi_n=\Psi_{n-1}b$ for $n\ge1$, and
$\Psi_0=\Phi_{\sigma_p}$. Consequently $\Psi$ induces a
linear map on algebraic Hochschild homology.
\end{proposition}

\begin{proof}
Let $\phi:\mathcal A\to\mathcal B$ be either shell
embedding. Apply $\phi^{\otimes n}$ to each term of
\eqref{eq:Hochschild-boundary}. Since
$\phi(a_ja_{j+1})=\phi(a_j)\phi(a_{j+1})$, the resulting
term is exactly the corresponding term of
$b\phi^{\otimes(n+1)}$. The cyclic last term obeys the
same identity. Thus both shell embeddings induce chain
maps, and their fixed linear combination does also.
Equation \eqref{eq:signed-embeddings} gives degree zero.
Chain maps send cycles to cycles and boundaries to
boundaries, so the homology map is well defined.
\end{proof}

Whether the induced map is nonzero on a primitive quotient
depends on that quotient. A comparison with the topological
or completed complex used in a trace formula is also needed
before this algebraic map can enter the arithmetic comparison.

The correction matters already in degree one. Restrict
to elementary chains whose group components are the identity,
and suppress the unchanged $\mathbb A_S$ coefficients.
The local factor of $\Psi_1$ is
\begin{equation}\label{eq:degree-one-factor}
 H_p=e_0\otimes e_0-p^{-1}e_1\otimes e_1.
\end{equation}
In contrast, tensoring the degree-zero shell map with
itself gives local factor
\begin{equation}\label{eq:naive-degree-one}
 \sigma_p\otimes\sigma_p
 =e_0\otimes e_0-p^{-1}(e_0\otimes e_1+e_1\otimes e_0)
      +p^{-2}e_1\otimes e_1.
\end{equation}
It has mixed-shell terms and a different outer-shell
coefficient. The proof of
Proposition~\ref{prop:corrected-chain-map} applies to
\eqref{eq:degree-one-factor}; multiplicativity of
$\Phi_{\sigma_p}$ cannot be used to justify the other
expression.

\begin{proposition}\label{prop:chain-Fourier-defect}
The degree-one local factor $H_p$ is not fixed by
componentwise local Fourier transformation. For every
$x,y\in B_{-1}$,
\begin{equation}\label{eq:chain-Fourier-defect}
 H_p(x,y)=0,\qquad
 ((\mathscr F_p\otimes\mathscr F_p)H_p)(x,y)
       =-\frac{(p-1)^3}{p^2}\ne0.
\end{equation}
\end{proposition}

\begin{proof}
The shells supporting $e_0,e_1$ do not meet $B_{-1}$,
so $H_p$ vanishes there. Formula \eqref{eq:ball-fourier}
gives
\begin{equation}\label{eq:shell-transform-pair}
 \mathscr F_pe_0=\ind_{B_0}-p^{-1}\ind_{B_1},\qquad
 \mathscr F_pe_1=p\ind_{B_{-1}}-\ind_{B_0}.
\end{equation}
Their values on $B_{-1}$ are $1-p^{-1}$ and $p-1$.
Therefore the transformed local tensor factor has value
\begin{equation}\label{eq:chain-Fourier-value}
 (1-p^{-1})^2-p^{-1}(p-1)^2
    =(p-1)^2(p^{-2}-p^{-1})
    =-(p-1)^3/p^2.
\end{equation}
The number is nonzero because $p>1$.
\end{proof}

This is an obstruction to the straightforward componentwise
Fourier symmetry of the corrected chain map. It does not
assert that Fourier transformation is itself a chain map
for the pointwise coefficient algebra. Indeed, Fourier
transformation exchanges pointwise multiplication and
additive convolution, so any geometric use of Fourier
duality must specify how the corresponding products are
related. The proposition isolates the local coefficient
failure without assuming such a relation.

The degree-zero local factor remains Fourier fixed:
$\mathscr F_p\sigma_p=\sigma_p$. Degree one already
requires additional structure. Thus the corrected chain
construction solves the extension problem at the level
of the algebraic boundary, while exposing a separate
duality problem. An eventual primitive comparison cannot
infer higher-degree Fourier compatibility from the
degree-zero identity.

\section{Two-prime chain compatibility}

The signed chain construction has a two-prime square at the
algebraic level. This square follows from local tensor embeddings,
not merely from the scalar Euler factors. It does not involve a
primitive projection or a Hermitian pairing.

For distinct primes $p,q\notin S$, write $e_{p,i}$ for the
$i$th local shell indicator at $p$, where $i\in\{0,1\}$,
and use $e_{q,j}$ at $q$. Set
\begin{equation}\label{eq:chain-coefficients}
 c_{p,0}=1,\qquad c_{p,1}=-p^{-1},\qquad
 c_{q,0}=1,\qquad c_{q,1}=-q^{-1}.
\end{equation}
Use a superscript $S$ to indicate the source place set of a
shell embedding or chain map. For example,
$\Psi_{p,n}^{S}$ maps $C_n(\mathcal A_S)$ to
$C_n(\mathcal A_{S\cup\{p\}})$.
All target local coordinates are placed in the same fixed
order when comparing compositions.

\begin{proposition}\label{prop:chain-square}
For every $n\ge0$,
\begin{equation}\label{eq:chain-square}
 \Psi_{q,n}^{S\cup\{p\}}\Psi_{p,n}^{S}
     =\Psi_{p,n}^{S\cup\{q\}}\Psi_{q,n}^{S}.
\end{equation}
Both compositions are chain maps. Their degree-zero
component is the coefficient map obtained by tensoring
with $\sigma_p\otimes\sigma_q$ and retaining the old
group element.
\end{proposition}

\begin{proof}
Each $\gamma\in\Gamma_{S\cup\{p\}}$ is a $q$-adic
unit, because $q\notin S\cup\{p\}$. Thus the
$q$-shell embeddings are defined on the full intermediate
algebra, including its new $p$ group elements. The analogous
statement holds with $p$ and $q$ exchanged.

On an elementary term $fU_\gamma$ from $\mathcal A_S$,
the composition of the $p$-shell $i$ embedding and
the $q$-shell $j$ embedding is
\begin{equation}\label{eq:two-shell-map}
 fU_\gamma\longmapsto(e_{p,i}\otimes e_{q,j}\otimes f)U_\gamma.
\end{equation}
Changing the order of the embeddings only interchanges the
two displayed local coordinates. The chosen common order
makes the maps equal. Call this algebra embedding
$\Phi_{i,j}$.

Composition of tensor powers equals the tensor power
of the composition. Expand the signed maps by linearity
to obtain, in either order,
\begin{equation}\label{eq:four-chain-terms}
 \sum_{i,j\in\{0,1\}}c_{p,i}c_{q,j}
            \Phi_{i,j}^{\otimes(n+1)}.
\end{equation}
This proves \eqref{eq:chain-square}. Each summand is a
chain map by the proof of
Proposition~\ref{prop:corrected-chain-map}; their combination is
therefore a chain map. At degree zero, factor the two
local sums in \eqref{eq:four-chain-terms} to recover
$\sigma_p\otimes\sigma_q$.
\end{proof}

No product of the signed maps was assumed to be
multiplicative in this proof. The calculation used the
four genuine algebra embeddings before combining their
induced chain maps. This distinction permits an exact
chain square despite
\cref{thm:local-correspondence:tensor-obstruction}.
The construction takes place over $\C$, where the
rational coefficients $p^{-1}$ and $q^{-1}$ are allowed.
It does not assert an integral correspondence with
integer multiplicities.

For a finite set $F$ of newly adjoined primes, the same
construction gives
\begin{equation}\label{eq:finite-chain-map}
 \Psi_{F,n}=
 \sum_{\mathbf i\in\{0,1\}^{F}}
 \left(\prod_{p\in F}c_{p,i_p}\right)
 \Phi_{\mathbf i}^{\otimes(n+1)},
\end{equation}
where $\Phi_{\mathbf i}$ tensors with the corresponding
product of shell indicators. Its degree-zero coefficient
is $\bigotimes_{p\in F}\sigma_p$. The formula is
independent of the order of adjoining the primes, since
each finite shell tensor product is independent of
that order after the canonical coordinate permutation.
Every summand is a chain map, so the formula gives a
finite algebraic directed system of chain maps.

The local Fourier failure also persists in all
positive degrees. The one-prime local factor in degree
$n$ is
\begin{equation}\label{eq:degree-n-factor}
 H_{p,n}=e_0^{\otimes(n+1)}-p^{-1}e_1^{\otimes(n+1)}.
\end{equation}
It vanishes on $B_{-1}^{n+1}$. By
\eqref{eq:shell-transform-pair}, its componentwise
Fourier transform is constant there with value
\begin{equation}\label{eq:degree-n-Fourier}
 (p-1)^{n+1}\bigl(p^{-(n+1)}-p^{-1}\bigr).
\end{equation}
This vanishes at $n=0$ and is strictly negative for
every $n\ge1$. Consequently the gap and Fourier
symmetry of the degree-zero tensor factor do not
extend to this componentwise construction in any
positive degree.

The algebraic boundary and the two-prime square are
respected by \eqref{eq:finite-chain-map}. The nonzero
local value in \eqref{eq:degree-n-Fourier} identifies
the componentwise Fourier symmetry that fails despite
those compatibilities.

\section{Error bounds for local truncations}

The convergent operator series give quantitative
approximations independent of a chosen test function.
Let
\begin{equation}\label{eq:KN}
 K_{p,\sigma}^{(N)}
   =L\sum_{k=1}^Np^{-k\sigma}(T_{kL}+T_{-kL}).
\end{equation}
The absolute sum of the omitted operator norms gives
\begin{equation}\label{eq:K-tail}
 \|K_{p,\sigma}-K_{p,\sigma}^{(N)}\|
       \le\frac{2Lp^{-(N+1)\sigma}}{1-p^{-\sigma}}.
\end{equation}
For arbitrary $f,g\in L^2(\R)$, the associated error in
the matrix coefficient is bounded by the right side
times $\|f\|_2\|g\|_2$. This estimate proves convergence
of the one-prime form on the full ambient Hilbert space.

For tests in $\T_a$, a different and stronger conclusion
holds once $NL\ge2a$: every omitted correlation value
vanishes, and the matrix-coefficient error is exactly
zero. A compact support argument supplies exactness;
the geometric bound supplies uniformity for functions
with no compact support. The two conclusions have
different domains.

The same comparison applies to the inverse. If
$R_N=\sum_{k=0}^Nr^kT_{kL}$, then
\begin{equation}\label{eq:inverse-residual}
 A_pR_N-I=-r^{N+1}T_{(N+1)L}.
\end{equation}
Its residual norm is exactly $r^{N+1}$, since the
translation is unitary. The inverse error is at most
$r^{N+1}/(1-r)$. At $p=2$, for example, $r=2^{-1/2}$;
the residual after $N$ is $2^{-(N+1)/2}$. A small
residual is a statement about this bounded local
inverse, not a bound for an arithmetic primitive
comparison.

The logarithm itself has the estimate
\begin{equation}\label{eq:log-tail}
 \left\|\log G_{p,\sigma}
 +\sum_{k=1}^N\frac{p^{-k\sigma}}k(T_{kL}+T_{-kL})\right\|
 \le\frac{2p^{-(N+1)\sigma}}{(N+1)(1-p^{-\sigma})}.
\end{equation}
It follows by replacing $1/k$ by $1/(N+1)$ in the
tail. Together with \eqref{eq:K-tail}, this gives a
direct local justification for differentiating after
uniform truncation on $\sigma\ge\sigma_0>0$.

For finite $F$, the matrix-coefficient errors from
truncating $\sum_{p\in F}K_{p,\sigma}$ are at most
the sum of their individual bounds. This elementary
estimate gives explicit error control for finite calculations with
different cutoffs for different primes. It supplies
no all-prime convergence statement at $\sigma=1/2$,
where summability over the prime index is a separate
question. Support-admissible tests avoid that question
for the distributional evaluation by making only
finitely many local coefficients relevant.

\section{Pairings and the geometric comparison}

The tensor obstruction occurs before a primitive pairing has been
specified. To obtain an intersection interpretation of
\eqref{eq:operator-local-Weil}, one needs an arithmetic complex
associated with $\mathcal A_S$, a primitive sector in its homology
or cohomology, and a pairing obtained from a specified duality and
trace. The analytic Sonin space would then require an actual
comparison map to that sector. Identifying vector spaces is only
one part of this construction.

An arithmetic comparison requires a primitive cohomological sector
$P_S$, a Hermitian pairing $I_S$ on that sector, and a complex-linear
test map $\mathcal G_S:\T\to P_S$. Their construction must establish
the finite comparison
\begin{equation}\label{eq:geometric-target}
 \begin{aligned}
 &I_{S\cup\{p\}}(\mathcal G_{S\cup\{p\}}f,\mathcal G_{S\cup\{p\}}g)
 -I_S(\mathcal G_Sf,\mathcal G_Sg)\\
 &\qquad=W_p(f*g^*)
 \end{aligned}
\end{equation}
after all degree and boundary terms have been accounted for.
The plus sign is forced by \eqref{eq:finite-increment} and the
desired identity $B_W=-I$.

The operator $K_{p,1/2}$ gives the right side a bounded analytic
realization for one prime. Defining $I_S$ by summing those
operators would transport the target arithmetic form. An
independently defined intersection pairing must instead be
compared with these operator matrix coefficients.

In the identity $\langle\theta_Sf,\eta_Sf\rangle=\|f\|^2$, the two
arguments belong to different images, and imposing one image as
the primitive sector would need a reason from the arithmetic
complex. Replacing $\eta$ by $\theta$ changes the form by
\eqref{eq:metric-defect}; differentiating that form changes it
again. These operations cannot be interchanged without changing
the proposed comparison.

At a two-prime stage, a primitive projection could introduce
mixed terms even though the analytic multipliers commute.
An equality of maps on ambient Hilbert spaces does not prove
compatibility of a pairing after those projections. The explicit
mixed translation term \eqref{eq:mixed-translations} gives a
test for any construction that claims to use inherited analytic
norms without modification.

Part II, Theorems 3.4 and 9.2 \cite{PartII}, rejects two additional explicit choices:
the ambient pointwise polar normalization on a nonzero entire
carrier, and an orthogonal nonpositive enlargement whose
increment is required to equal a positive local witness.
Those obstructions leave mixed pairings, boundary corrections,
and changing primitive projections unresolved. Together with
\cref{thm:local-correspondence:tensor-obstruction}, they narrow
the construction problem to correspondences with additional
structure beyond a tensor algebra map and a normalized Hilbert
inclusion.

\section{Completion and dilation domains}

For a finite set $F$ of primes, the product
$A_F=\prod_{p\in F}A_p$ is bounded and invertible. Repeated use
of \eqref{eq:Ap-bounds} gives
\begin{equation}\label{eq:finite-product-bounds}
 \prod_{p\in F}(1-p^{-1/2})\|u\|_2
 \le\|A_Fu\|_2
 \le\prod_{p\in F}(1+p^{-1/2})\|u\|_2.
\end{equation}
These are finite bounds. They do not provide uniform control
as $F$ grows through all primes. Neither the operator-norm
convergence of each one-prime inverse nor the compact-support
stabilization of the arithmetic distribution establishes
convergence of this infinite product on the critical line.

The local derivative behaves better with respect to a finite
product because of \eqref{eq:log-additivity}. For
$G_{F,\sigma}=\prod_{p\in F}G_{p,\sigma}$,
\begin{equation}\label{eq:finite-log-sum}
 \frac{d}{d\sigma}\log G_{F,\sigma}
       =\sum_{p\in F}K_{p,\sigma}.
\end{equation}
At $\sigma=1/2$, pairing with a fixed compact test correlation
eventually stabilizes to the finite-prime part of the explicit
formula. This is a statement about matrix coefficients on the
test algebra, not operator-norm convergence on the ambient
Hilbert space. Its distinction from an infinite product remains
even though both use the same local factors.

The fixed Sonin gap also places a constraint on dynamics.
For $b>0$, let $D_bv(x)=b^{1/2}v(bx)$ on $L^2(\R)$.
If $v$ has spatial and Fourier gaps both equal to $\lambda$,
then $D_bv$ has the two guaranteed gaps
\begin{equation}\label{eq:dilation-gaps}
 \lambda/b\quad\hbox{and}\quad b\lambda,
\end{equation}
respectively. The second follows from
$\mathscr F(D_bv)(\xi)=b^{-1/2}(\mathscr Fv)(\xi/b)$.
Thus the definition of a fixed Sonin space does not by itself
supply an invariant representation of all dilations on that
space. A self-adjoint generator for the ambient scaling action
cannot simply be restricted without a domain argument.

The same issue arises for a proposed weight-one generator
$\Theta$. If a densely defined operator on a Hilbert space
satisfies $\Theta^*=I-\Theta$ with equality of domains, then
$D=(\Theta-\tfrac12I)/i$ is self-adjoint. Indeed,
$D^*=i(\Theta^*-\tfrac12I)=D$, and the assumed domain equality
is exactly what makes this an equality of operators. A formal
identity on a common core is insufficient unless the closures
and their domains are controlled.

An arithmetic spectral realization would also have to account
for every nontrivial zero with its multiplicity. The local
operators in this paper were built without using a preselected
zero list. That feature avoids circularity, but it does not
identify their spectra with the zeta zeros. The local
translation spectrum and a complete arithmetic zero
realization are different objects.

\section{Formal scope}

The companion source archive \cite{APCcode} contains the Lean 4
development. Its toolchain is \texttt{v4.30.0-rc2}; the Lake
manifest pins mathlib to commit
\nolinkurl{0f9072dd907c6e2e4264ab241a049cab50137f7c}.
The September 2026 source hashes are recorded in
\nolinkurl{research/lean-source-sha256.txt}, and
\nolinkurl{research/lean-results.md} states the proof scope.

The module \nolinkurl{lean/APC/SoninShell.lean} verifies the
shell-coefficient algebra underlying
Theorem~\ref{thm:local-correspondence:tensor-obstruction}.
The function \nolinkurl{APC.soninShellCoefficient} assigns the
complex values $1$, $-1/p$, and $0$ to integer shell indices
$0$, $1$, and all other indices. The theorem
\nolinkurl{soninShell_no_nonzero_scalar_idempotent}, in namespace
\texttt{APC}, proves for every natural $p>1$ and every
nonzero complex $c$ that these scaled values cannot all
satisfy $(c\sigma(n))^2=c\sigma(n)$.

The theorem \nolinkurl{soninShell_tensor_not_multiplicative}
then rejects multiplicativity of the scalar-to-shell
function map. These statements formalize the algebraic
contradiction at the two shell values. The development
does not define $p$-adic functions or the crossed-product
algebra, and it does not prove the uniqueness or Fourier
invariance of the local Sonin vector. Those passages
are the mathematical arguments of Sections 4 and 13.

In \nolinkurl{lean/APC/LocalFactor.lean}, the squared modulus
is computed by \nolinkurl{localWeight_formula}. Its bounds
are \nolinkurl{localWeight_lower} and \nolinkurl{localWeight_upper}.
The declarations
\nolinkurl{localWeight_defect_negative} and
\nolinkurl{localWeight_defect_positive} prove its strict
signs at the specified frequencies. The theorem
\nolinkurl{localPrimePowerSum_eq_zero_of_support}
proves support vanishing for the finite local sum.
All these declarations are in namespace \texttt{APC}.
The operator logarithm, its norm-convergent derivative,
and Fourier inversion in Theorem~\ref{thm:local-correspondence:logarithmic}
remain proofs in the manuscript. No Lean theorem is
claimed for the Hochschild maps or a primitive geometric
pairing. The distinction records exactly which parts
of the construction the formal statements cover.

\section{Conclusion}

The shell construction gives a bounded prime-adjoining
transition with exact Fourier duality, and its metric logarithm
has an operator derivative equal to the full local prime-power
distribution. The inherited norm defect has a different
kernel and nonzero mixed-prime terms. Support-admissible
windows recover the complete arithmetic distribution without
requiring an infinite Euler product.

The direct tensor map from the local Sonin vector fails
multiplicativity in the semilocal crossed-product algebra.
This failure holds for every nonzero scalar normalization and
therefore rejects that specific geometric upgrade of the
analytic transition. Signed shell embeddings give a corrected
algebraic Hochschild map and commuting prime transitions,
but componentwise Fourier symmetry fails in every positive
degree. A replacement duality must recover the local
logarithmic term. The primitive sign and the complete
arithmetic comparison remain unproved.

\begin{thebibliography}{9}
\small
\bibitem{CCM}
Alain Connes, Caterina Consani and Henri Moscovici.
\emph{Zeta zeros and prolate wave operators: Semilocal adelic operators}.
arXiv:2310.18423v2, 2024.
\url{https://arxiv.org/html/2310.18423v2}.

\bibitem{Connes2026}
Alain Connes.
\emph{The Riemann Hypothesis: Past, Present and a Letter Through Time}.
arXiv:2602.04022v1, 2026.
\url{https://arxiv.org/html/2602.04022v1}.

\bibitem{CCMar}
Alain Connes, Caterina Consani and Matilde Marcolli.
\emph{The Weil proof and the geometry of the adeles class space}.
arXiv:math/0703392v1, 2007.
\url{https://arxiv.org/html/math/0703392v1}.

\bibitem{PartII}
Matthew Long and The YonedaAI Collaboration.
\emph{Obstructions to Adelic Polarization}.
Part II of \emph{Adelic Polarization Conjecture}, companion manuscript,
September 2026.

\bibitem{CC2020}
Alain Connes and Caterina Consani.
\emph{Weil positivity and Trace formula, the archimedean place}.
arXiv:2006.13771v1, 2020.
\url{https://arxiv.org/html/2006.13771v1}.

\bibitem{APCcode}
Matthew Long and The YonedaAI Collaboration.
\emph{Adelic Polarization Conjecture: companion source archive}.
September 2026. Lean sources, toolchain, dependency manifest
and source hashes accompany the manuscripts.
\url{https://github.com/YonedaAI/adelic-polarization-conjecture}.
\end{thebibliography}
\end{document}
